% gwebmac.tex -- plain-TeX macros for documents produced by gweave.
% This is the GWEB analogue of CWEB's cwebmac.tex; it is independent macro code
% written to give gweave output the same page layout that cweave produces.
% Process a woven file with a plain-TeX engine, e.g. pdftex foo.tex.

% ---- fonts ------------------------------------------------------------------
% The font set mirrors cwebmac.tex so woven output matches cweave's typefaces.
\font\eightrm=cmr8                       % small roman for notes, headers, index
\font\ninerm=cmr9 \let\mc=\ninerm        % medium caps (\CEE/, \GO/, ...)
\let\sc=\eightrm                         % smallish caps
\font\titlefont=cmr7 scaled\magstep4     % big title on the contents page
\font\ttitlefont=cmtt10 scaled\magstep2  % typewriter title (for use in limbo)
\font\tentex=cmtex10                      % TeX extended set, for strings (as cweb)
\fontdimen7\tentex=0pt                    % no extra space after sentences
\let\mainfont=\tenrm
\let\cmntfont=\tenrm                      % font for comments (cf. cweb's \cmntfont)

% Language and prose helpers borrowed from cwebmac, for use in commentary; each
% is invoked with a trailing slash, e.g. ``written in \GO/''.
\def\CEE/{{\mc C\spacefactor1000}}
\def\GO/{{\mc Go\spacefactor1000}}
\def\UNIX/{{\mc U\kern-.05emNIX\spacefactor1000}}
\def\TEX/{\TeX}
% \\{id} sets an identifier in italics in prose, as in cweb. (The GWEB idiom is
% to write inline code as |id|; this is provided for porting cweb commentary.
% cweb's \& and \| are deliberately not redefined: \& is used here as a literal
% ampersand by the comment/section-name escaper, and gweave rewrites \| to a bar.)
\def\\#1{\leavevmode\hbox{\it#1\/\kern.05em}}
% A thin space that works in either math or text mode (cwebmac's \,).
\def\,{\relax\ifmmode\mskip\thinmuskip\else\thinspace\fi}

% \.{text} sets verbatim/typewriter text in prose (as in cwebmac). Inside it the
% escaped specials \\ \{ \} \& \# \% \$ \^ \_ \~ print as the literal characters,
% so e.g. \.{@\^} prints the control code @^. Used for file names, command names,
% and control codes in commentary and section titles. (Plain TeX's \. is a dot
% accent; we override it -- gweave documents need \tt strings, not accents.)
% This is a plain \def (not e-TeX's \protected, so bare `tex' works for the DVI
% route): the running-head and contents machinery below is written to carry the
% title without ever expanding \. (see \setmark and \runhead).
\def\.#1{{\tentex
  \def\\{\char92 }\def\{{\char123 }\def\}{\char125 }%
  \def\&{\char38 }\def\^{\char94 }\def\_{\char95 }\def\~{\char126 }%
  \def\${\char36 }\def\#{\char35 }\def\%{\char37 }%
  #1}}

% ---- page layout ------------------------------------------------------------
\newdimen\pagewidth  \pagewidth=6.5in
\newdimen\pageheight \pageheight=8.7in
\newdimen\colwidth   \colwidth=.5\pagewidth \advance\colwidth by-10pt
\hsize=\pagewidth
\vsize=\pageheight
\parindent=1em
\parskip=0pt
\hbadness=10000
\tolerance=2000
\newdimen\indentunit \indentunit=1.4em % one indentation step in displayed code
\newcount\tmpcount \newdimen\tmpdimen  % scratch registers (avoid e-TeX exprs)
\newskip\secskip     \secskip=12pt plus 4pt minus 3pt % space between sections
\newcount\codepenalty \codepenalty=9999 % discourage page breaks inside code
                                          % (high but finite, so >1-page code can still break)
\newdimen\commentkern \commentkern=-1.3pt % tighten the "//" comment marker
\newcount\secpagedepth \secpagedepth=1 % starred sections this deep start a page

% "++" and "--" as tight, small, raised symbols, exactly as cweave sets them:
% seven-point signs kerned together so the pair reads as one operator, not two.
\newbox\PPbox % the "++" symbol
\setbox\PPbox=\hbox{\kern.5pt\raise1pt\hbox{\sevenrm+\kern-1pt+}\kern.5pt}
\def\PP{\copy\PPbox}
\newbox\MMbox % the "--" symbol
\setbox\MMbox=\hbox{\kern.5pt\raise1pt\hbox{\sevensy\char0\kern-1pt\char0}\kern.5pt}
\def\MM{\copy\MMbox}

% Go's short variable declaration ":=", as one symbol you may redefine. WEB set
% Pascal's ":=" as a left arrow through \K; \K is the same hook for Go, so a
% document that prefers an arrow need only say, in its limbo,
%     \let\K=\Leftarrow
\def\K{{:}{=}}

% Go's bit-clear operator "&^", likewise redefinable. It has no C counterpart, so
% its glyph is a choice rather than a tradition: a circled slash, read as "and
% not". Say, for instance, \let\AN=\setminus in your limbo to prefer another.
\def\AN{\oslash}

% ---- running headers --------------------------------------------------------
% The document title is the upper-cased job name; the current group title is set
% by each starred section. Each section deposits a \mark of the form
% {section-number}{group-title}; the headline reads it back through \topmark.
% Upper-case the (expanded) job name for the running header and contents title.
% \title is shown, as written, in the running header and on the contents page. It
% defaults to the upper-cased job name -- so a web with no \def\title still gets
% an all-caps header, as in cweb -- but a \def\title{...} in limbo is shown
% verbatim, mixed case and all: cweb likewise upper-cases only the default
% job-name title, never one the author supplies.
\edef\title{\jobname}
\def\titleaux#1{\uppercase{\gdef\title{#1}}}
\expandafter\titleaux\expandafter{\title}
\def\Gtitle{\title}
% A snapshot of the default (job-name) title, taken before limbo runs, so a
% localization can tell whether the author supplied a \title of their own:
% \ifx\title\defaulttitle is true exactly when none was given.
\let\defaulttitle\title
% The current group (chapter) title, upper-cased, is kept in a token register so
% each \mark can capture a *snapshot* of it (via \the) rather than a reference to
% a macro -- otherwise every page would show whatever group is current at shipout
% time, i.e. the next chapter. \the inserts the tokens without expanding them, so
% a \. in the title survives into \topmark for typesetting in the header.
\newtoks\grouptoks
\def\setmark#1{\mark{{#1}{\the\grouptoks}}}
\def\takeone#1#2{#1}
\def\taketwo#1#2{#2}
\def\empty{}

\newif\iftitle % true only while shipping the contents page (suppresses the header)
\def\contentspagenumber{0} % page number of the contents page; redefine in limbo
\def\sechead{\S\thinspace\expandafter\takeone\tm} % left/right end: section number
\def\grouphead{\expandafter\taketwo\tm}   % middle: current group title
% Font for the running-head title and group title (small roman, as in cweb).
% A localization file may rebind it -- e.g. kotexgweb.tex pairs it with a sans
% hangul font so Korean in the header matches the upper-cased Latin.
\def\headfont{\eightrm}
% The running head sets the document title in \Gtitlefont and the group (chapter)
% title in \groupfont; both default to \headfont, but a localization may split
% them -- e.g. kotexgweb keeps the group title sans (to match the upper-cased
% Latin) while setting an author-supplied document title in serif, to match its
% serif rendering on the contents page.
\def\Gtitlefont{\headfont}
\def\groupfont{\headfont}
% Font for the small print: cross-reference notes, the contents columns, and the
% usage note in the section-name list. A localization file may rebind it -- e.g.
% kotexgweb.tex pairs it with an 8pt hangul font so Korean there is small too.
\def\smallfont{\eightrm}
% \runhead builds the header as one \pagewidth-wide box (used both by the
% normal \headline and by the index's own output routine).
\def\runhead{%
  % Capture the mark tokens without expanding them (a \. in the title must not be
  % expanded here); \expandafter expands \topmark exactly once into \tm.
  \expandafter\def\expandafter\tm\expandafter{\topmark}%
  \ifx\tm\empty\expandafter\def\expandafter\tm\expandafter{\firstmark}\fi
  \ifx\tm\empty\def\tm{{}{\Gtitle}}\fi
  \ifodd\pageno
    \hbox to\pagewidth{\mainfont\sechead\enspace{\Gtitlefont\Gtitle}\hfil
      {\groupfont\grouphead}\enspace\folio}%
  \else
    \hbox to\pagewidth{\mainfont\folio\enspace{\groupfont\grouphead}\hfil
      {\Gtitlefont\Gtitle}\enspace\sechead}%
  \fi}
\headline={\iftitle\hfil\else\runhead\fi}
\footline={\hfil}

% ---- hyperlinks --------------------------------------------------------------
% Each section deposits a numbered destination; section numbers shown as
% references, in the index, in cross-reference notes, and in the contents link
% to those destinations, drawn in blue like cweave. Two back ends are supported:
%   * pdfTeX/luaTeX primitives, when producing PDF directly (\pdfoutput>0);
%   * \special{pdf:...} commands for the  tex -> dvipdfmx  route. As in cweb,
%     that route is requested by defining \pdf, e.g.
%         tex "\let\pdf+ \input foo.tex"   then   dvipdfmx foo.dvi
% With a plain dvi engine and no \pdf, the links and bookmarks are omitted.
% For the dvipdfmx route, destinations are named with roman numerals (as in
% cwebmac); links and bookmarks reference those names, so they work even though
% gweave emits all the bookmarks together at the end of the document.
\newif\ifpdf      % native pdfTeX/luaTeX PDF output
\newif\ifdvipdf   % dvipdfmx \special{pdf:...} route
\ifx\pdfoutput\undefined\else\ifnum\pdfoutput>0 \pdftrue\fi\fi
\ifpdf\else\ifx\pdf\undefined\else\dvipdftrue\fi\fi
\newbox\linkbox
\def\dest#1{%
  \ifpdf \pdfdest num #1 fith\relax
  \else\ifdvipdf
    \smash{\raise\baselineskip\hbox to0pt{%
      \special{pdf: dest (\romannumeral#1) [ @thispage /FitH @ypos ]}\hss}}%
  \fi\fi}
% The destination for the `Names of the sections' page. It is emitted only when
% that list is (\ifsecs), so \nosecs leaves no stray whatsit---and hence no
% blank page---on the sheet the index ejects into.
\def\secdest#1{\ifsecs\dest{#1}\fi}
\def\link#1#2{%
  \ifpdf
    \pdfstartlink attr{/Border[0 0 0]}goto num #1\relax
    \pdfliteral{0 0 1 rg}#2\pdfliteral{0 g}\pdfendlink
  \else\ifdvipdf
    \setbox\linkbox=\hbox{\special{pdf: bc [ 0 0 1 ]}#2\special{pdf: ec}}%
    \special{pdf: ann width \the\wd\linkbox\space height \the\ht\linkbox\space
       depth \the\dp\linkbox\space << /Subtype /Link /Border [ 0 0 0 ]
       /A << /S /GoTo /D (\romannumeral#1) >> >>}\box\linkbox\relax
  \else #2\fi\fi}
\def\s#1{\link{#1}{#1}}                        % a linked section number
\def\sD#1{\link{#1}{$\underline{\hbox{#1}}$}}  % linked + underlined (definition)
% A PDF outline (bookmark) entry: #1 group depth, #2 section number, #3 number of
% direct children (pdfTeX's open/close model), #4 the plain-text title. The
% dvipdfmx route nests by depth and links to the named destination.
\def\bookmark#1#2#3#4{%
  \ifpdf \pdfoutline goto num #2 count #3 {\pdfescapestring{#4}}%
  \else\ifdvipdf
    \special{pdf: outline #1 << /Title (#4)
       /A << /S /GoTo /D (\romannumeral#2) >> >>}%
  \fi\fi}
% \pdfURL{text}{url}: `text' set in blue as a clickable link to `url', as in
% cwebmac; with no PDF back end both are printed, the URL in parentheses and
% typewriter. In the url write @@ for an at-sign (the usual .w escape),
% \TILDE/ for ~, and \UNDER/ for _, as in cwebmac.
\def\pdfURL#1#2{%
  \ifpdf
    \pdfstartlink attr{/Border[0 0 0]}
      user{/Subtype /Link /A << /S /URI /URI (#2) >>}\relax
    \pdfliteral{0 0 1 rg}#1\pdfliteral{0 g}\pdfendlink
  \else\ifdvipdf
    \setbox\linkbox=\hbox{\special{pdf: bc [ 0 0 1 ]}#1\special{pdf: ec}}%
    \special{pdf: ann width \the\wd\linkbox\space height \the\ht\linkbox\space
       depth \the\dp\linkbox\space << /Subtype /Link /Border [ 0 0 0 ]
       /A << /S /URI /URI (#2) >> >>}\box\linkbox\relax
  \else #1 ({\tt#2})\fi\fi}
{\catcode`\~=12 \gdef\TILDE/{~}} % ~ in a URL
{\catcode`\_=12 \gdef\UNDER/{_}} % _ in a URL
% \runheadline places \runhead like plain TeX's \makeheadline, for use inside
% the index's custom \shipout (where the plain \headline machinery is bypassed).
% It is deliberately not named \makeheadline: plain TeX's own \makeheadline is
% live here (\output is \plainoutput on ordinary pages), and shadowing it would
% make the header-suppressing \headline logic yield to an unconditional \runhead
% on the contents page, where \iftitle is true.
\def\runheadline{\vbox to 0pt{\vskip-22.5pt
  \hbox to\pagewidth{\vbox to8.5pt{}\runhead}\vss}\nointerlineskip}

% A starred section records itself in \contentsfile as it is shipped, so \con
% can build the contents in one pass; the page number is captured by the deferred
% \write, the title is stashed in a control sequence. \contentsfile (the file the
% data goes to) and \readcontents (how \con reads it back) are redefinable in
% limbo, as in cwebmac: to split a long document into separately-typeset pieces
% and gather their contents in the last one, give each piece its own
% \contentsfile and let the last \readcontents \input them all. The file is
% opened lazily, by \opentoc at the first starred section, so that a limbo
% redefinition of \contentsfile is already in force when it opens.
\def\contentsfile{\jobname.toc}
\def\readcontents{\input \contentsfile }
\newwrite\toc
\newtoks\toctoks % holds a starred section's title, unexpanded, for the deferred \write
\newif\iftocopen
\def\opentoc{\iftocopen\else\tocopentrue\immediate\openout\toc=\contentsfile \fi}

% ---- section openers --------------------------------------------------------
% \M{n} : an ordinary (unstarred) section, run-in bold number. The
% \vfil\penalty-100\vfilneg makes the section boundary the preferred page break
% (the \vfil fills the page there, so a short page costs nothing): as in cweave,
% TeX breaks between sections rather than inside one, so a whole section -- its
% commentary and its code -- stays together on a page when it fits.
\def\M#1{\par\emitdate\vfil\penalty-100\vfilneg\vskip\secskip
  \setmark{#1}\noindent\dest{#1}{\bf #1.\quad}\rightskip=0pt\ignorespaces}

% \N{depth}{n}{title} : a starred section. Major sections (depth < secpagedepth)
% begin a new page; the first one does not, so it follows any limbo banner.
\newif\iffirst \firsttrue
\def\N#1#2#3{\par\emitdate
  % Snapshot the upper-cased title into \grouptoks; \uppercase only upper-cases
  % the letter tokens, leaving control sequences (\.) intact for the header.
  {\uppercase{\global\grouptoks={#3}}}%
  \setmark{#2}% set the mark before any page break, so the new page's \topmark
  \iffirst \firstfalse \vskip 6pt
  \else \ifnum#1<\secpagedepth \vfil\eject
        \else \penalty-200\vfil\penalty-200\vfilneg\vskip\secskip \fi
  \fi
  \opentoc
  % The title (#3) must reach the contents run unexpanded. \unexpanded is an e-TeX
  % primitive, absent from Knuth's tex, so we use plain TeX's own device instead: a
  % token register, whose contents \the emits without re-expanding. The register is
  % set after the page break above and read by the deferred \write at ship-out; a
  % starred section always begins its own page (\secpagedepth), so no later section
  % overwrites it before then. Only \the\pageno is left for the \write to fill in.
  \toctoks={#3}%
  \write\toc{\string\Tline{#1}{#2}{\the\toctoks}{\the\pageno}}%
  \noindent\dest{#2}{\bf #2.\quad #3.\quad}\rightskip=0pt\ignorespaces}

% ---- \datethis: timestamp the listing just above section 1 ------------------
% cweave's \datethis (said in limbo) prints the date and time, in \eightrm (its
% \sc), on a left-flush line above the first section. \today and \hours compute
% them; \datefont and \dateat -- the font and the date/time separator -- are
% overridable, so the Korean backend can localize the wording and the font size.
\def\today{\ifcase\month\or
  January\or February\or March\or April\or May\or June\or
  July\or August\or September\or October\or November\or December\fi
  \space\number\day, \number\year}
\newcount\twodigits
\def\hours{\twodigits=\time \divide\twodigits by60 \printtwo
  \multiply\twodigits by-60 \advance\twodigits by\time :\printtwo}
\def\gobbleone1{}
\def\printtwo{\advance\twodigits by100
  \expandafter\gobbleone\number\twodigits \advance\twodigits by-100 }
\def\datefont{\eightrm}  % cweb sets this line in \sc, i.e. \eightrm (8pt)
\def\dateat{\ at\ }      % between the date and the time
\def\dateline{\leftline{\datefont\today\dateat\hours}\bigskip}
\let\datehook\empty
\def\datethis{\let\datehook\dateline}
\def\emitdate{\datehook\global\let\datehook\empty} % fires once, at section 1

% ---- displayed Go code ------------------------------------------------------
\newdimen\codeindent \codeindent=\indentunit % displayed code is indented this
                                                % far (about one tab) from the margin
% \penalty\codepenalty before the code keeps it joined to its commentary (the
% section breaks as a unit); like the penalty inside \GL it is high but finite,
% so a section too tall for a page can still break, preferring the commentary.
\def\B{\par\penalty\codepenalty\vskip 3pt\begingroup
  \rightskip=0pt plus 1fil \parfillskip=0pt plus 1fil
  \hfuzz=2pt \emergencystretch=3em \parindent=0pt}
% \Br: run-in code start for a code-only section with no commentary -- no \par,
% so the first code line (set with \Lr) follows the section number; used by
% cweave-style code-only sections that have no name and no commentary.
\def\Br{\penalty\codepenalty\begingroup
  \rightskip=0pt plus 1fil \parfillskip=0pt plus 1fil
  \hfuzz=2pt \emergencystretch=3em \parindent=0pt}
\def\E{\par\endgroup\vskip 3pt}
% One displayed code line: #1 indentation steps, #2 the body (math chunks $...$
% joined by \GS). The whole display is offset by \codeindent from the margin;
% long lines wrap at a \GS, continuation indented one step deeper.
% \penalty\codepenalty before each line discourages (but, being finite, does not
% forbid) a page break between consecutive code lines, so a code part prefers to
% stay together; a code part longer than a page can still break.
\def\GL#1#2{\par\penalty\codepenalty
  \tmpcount=#1 \tmpdimen=\tmpcount\indentunit \advance\tmpdimen by\codeindent
  \hangindent=\tmpdimen \advance\hangindent by\indentunit \hangafter=1
  \noindent\kern\tmpdimen \strut #2}
% \Lr: the run-in first code line of a code-only section -- no \par and no
% indent, so it continues on the section-number line; later lines use \GL.
\def\Lr#1#2{\strut #2}
\def\GS{\penalty0\hskip .222em plus .12em minus .08em\relax} % breakable code space (cweb medmuskip)
\def\punct{\penalty0\hskip .1667em plus .09em minus .06em\relax} % after a comma (cweb thinmuskip, 3mu)
\def\rel{\penalty0\hskip .2778em plus .15em minus .10em\relax} % around a relation (cweb thickmuskip, 5mu)
\def\W{\penalty0\hskip .333em plus .167em minus .111em\relax} % between two words (cweb's text interword space)
\def\BS{\penalty0\hskip .5em plus .2em minus .12em\relax} % wider structural space around a block brace (cweb's \5)
\def\CS{\penalty0\hskip 1em plus .3em minus .1em\relax}  % generous space before a comment (cweb's \5\5\quad)
\def\SO{\penalty0\hskip 0pt\relax}                       % @| optional break
\def\BL{\par\penalty\codepenalty\smallskip}             % @# blank line in code
\def\BK{\par\penalty\codepenalty\vskip 2.5pt} % a blank source line: a little air

% A |...| fragment from the TeX parts (prose, comments, section names) is wrapped
% in \PB, cweb's ``program brackets'': the code material is math (it uses
% \mathord and friends), so \PB supplies the enclosing $...$ only when TeX is
% not already in math mode. That way a |...| the author put inside a $...$ span,
% or inside a text-mode \halign cell of a $$...$$ display, is set correctly
% either way -- TeX decides the mode, so gweave need not guess it.
\def\PB#1{\ifmmode#1\else$#1$\fi}
% ---- code token fonts (called inside the math of \GL) -----------------------
\def\ID#1{\hbox{\it #1\/}}      % identifier
\def\KW#1{\hbox{\bf #1}}        % reserved word or predeclared type
\def\ST#1{\hbox{\tentex #1}}    % string, rune, or verbatim text (cweb's \tentex)
\def\MAC#1{\hbox{\tentex #1}}   % const name, set typewriter like cweb's @d macro
\def\SP{{\tt\char32 }}          % a visible space inside a string (cweb's \SP)
\def\nil{\mathord{\Lambda}}     % Go's nil, shown as cweb shows C's NULL (\Lambda)
\def\CM#1{\hbox{\cmntfont #1}}  % comment
\def\NU#1{\hbox{\rm #1}}        % decimal/floating numeric literal
% Non-decimal literals, as in cweb: octal with a raised circle and oldstyle
% italic digits, hex with a superscript "#" in typewriter, binary with a "b".
\def\oct#1{\hbox{$^\circ$\kern-.1em{\it #1\/}}}
\def\hex#1{\hbox{$^{\scriptscriptstyle\#}$\tt #1}}
\def\bin#1{\hbox{$^{\scriptscriptstyle b}$\tt #1}}

% ---- named-section references and headlines ---------------------------------
\def\xsec#1#2{$\langle\,$#1\hbox{\eightrm\kern.4em\s{#2}}$\,\rangle$}
\def\X#1#2{\hbox{\xsec{#2}{#1}}}        % inline reference in code
\def\D#1#2{\par\vskip 3pt\noindent\hbox{\xsec{#2}{#1}${}\equiv{}$}\hskip 3pt}
\def\Dp#1#2{\par\vskip 3pt\noindent
  \hbox{\xsec{#2}{#1}${}\mathrel{+}\equiv{}$}\hskip 3pt}
% Run-in variants \Dr/\Dpr: used when a section has no commentary, so the
% definition's name and equivalence sign follow the section number on the same
% line (no \par), as cweave does for a code-only section.
\def\Dr#1#2{\hbox{\xsec{#2}{#1}${}\equiv{}$}\hskip 3pt}
\def\Dpr#1#2{\hbox{\xsec{#2}{#1}${}\mathrel{+}\equiv{}$}\hskip 3pt}

% ---- cross-reference notes printed under a definition -----------------------
\def\note#1{\par\penalty-1\vskip 2pt\noindent{\smallfont #1.}\par}
\def\U#1{\note{This code is used in section #1}}
\def\Us#1{\note{This code is used in sections #1}}
\def\A#1{\note{This code is also defined in section #1}}
\def\As#1{\note{This code is also defined in sections #1}}
% Other human-language strings: the running head over the section-name list and
% the two column headings on the contents page. Redefine these, and the four
% notes just above, in a localization file (e.g. kotexgweb.tex for Korean) to
% translate gweave's output; gweave itself needs no change.
\def\namesheading{NAMES OF THE SECTIONS}
\def\outsecname{Names of the sections} % title of the section-names PDF bookmark
\def\sectionword{Section}
\def\pageword{Page}
% The usage note printed beside an entry in the section-name list.
\def\Nused#1{Used in section #1}
\def\Nuseds#1{Used in sections #1}

% ---- selectively suppressing the back matter (as in cwebmac) ----------------
% Put any of these in limbo to drop part of the tail. Following cwebmac they are
% cumulative in reach: \noinx drops the index, the `Names of the sections' list,
% and the contents; \nosecs drops the list and the contents; \nocon drops the
% contents. The three flags guard the matching parts of \inx/\II, \fin/\NS,
% and \con below, so every group stays balanced whatever is switched off.
% \noatl is accepted for portability with cwebmac webs but does nothing: GWEB
% has no \.{@l} transliteration (Go source is Unicode), so there is no such note.
\newif\ifinx  \inxtrue   % include the index?
\newif\ifsecs \secstrue  % include the `Names of the sections' list?
\newif\ifcon  \contrue   % include the table of contents?
\def\noinx{\inxfalse\secsfalse\confalse}
\def\nosecs{\secsfalse\confalse}
\def\nocon{\confalse}
\def\noatl{}

% ---- index: title-less, set in two columns ----------------------------------
% The index is set in two columns. \inx stashes the preceding material (the
% user's `@* Index.' heading and prose) in \sbox, then narrows \hsize to one
% column and installs \coloutput: it fills the left column to the column height,
% saves it, fills the right column, and ships a page carrying the stashed
% material above the two columns. \fin flushes the last column.
\newbox\sbox \newbox\lbox
\def\lr{L}
\def\coloutput{%
  \if L\lr
    \global\setbox\lbox=\box255 \gdef\lr{R}%
  \else
    \shipout\vbox{\runheadline
      \vbox to\pageheight{\boxmaxdepth=\maxdimen
        \box\sbox\vss
        \hbox to\pagewidth{\box\lbox\hfil\box255}}}%
    \global\advance\pageno by1
    \global\setbox\sbox=\vbox{}\global\vsize=\pageheight \gdef\lr{L}%
  \fi}
\def\inx{\ifinx
  \par
  \vskip 9pt plus 2pt           % gap between the Index heading/prose and entries
  \output={\global\setbox\sbox=\box255}\eject % stash heading+prose, hold page
  \setbox\sbox=\vbox{\unvbox\sbox}%
  \global\vsize=\pageheight
  \global\advance\vsize by-\ht\sbox \global\advance\vsize by-\dp\sbox
  \global\hsize=\colwidth
  \global\output={\coloutput}\gdef\lr{L}%
  \begingroup
  \baselineskip=12pt \parindent=0pt \parskip=1pt plus .5pt \rm
  \rightskip=0pt plus 2.5em \tolerance=10000 \hyphenpenalty=10000
  \fi}
\def\II#1#2{\ifinx\noindent\hangindent2em #1:\kern.5em #2.\par\fi} % index entry
\def\IR#1{{\rm #1}}                          % @^...@>  roman index entry
\def\IT#1{{\tentex #1}}                      % @....@>  typewriter index entry
\def\9#1{{\rm #1}}                            % default for @:...@> entries
\def\IC#1{\9{#1}}

% ---- list of section names (a `Names of the Sections' page) -----------------
% \fin ends the index (\endindex) and starts the section-name list
% (\beginsecs); the two halves are split out so \nosecs can end the index and
% stop, and each is guarded by its own flag so the groups stay balanced.
\def\endindex{\par\vfil\eject  % seal/ship the current index column
  \if L\lr\else\null\vfil\eject\fi % a left column is pending: ship its page
  \endgroup                     % close the index paragraph group (from \inx)
  \global\output={\plainoutput}% back to ordinary single-column pages
  \global\hsize=\pagewidth \global\vsize=\pageheight}
\def\beginsecs{\vfil\eject
  \global\def\sechead{}% no section number in the header here
  \global\def\grouphead{\namesheading}%
  \begingroup
  \parindent=0pt \parskip=2pt plus 1pt \rightskip=0pt plus 2em \rm}
\def\fin{\ifinx\endindex\fi \ifsecs\beginsecs\fi}
\def\NS#1#2#3{\ifsecs\noindent\hangindent2em \xsec{#1}{#2}%
  \ifx\empty#3\else\quad{\smallfont#3}\fi\par\fi}

% ---- contents page: centered title, Section/Page columns, dot leaders --------
% \Tline{depth}{secno}{title}{page}: one contents entry. The title travels in
% the toc file itself (written unexpanded by \N), so the file is self-contained
% and \readcontents can gather several pieces' contents in one run.
\def\Tline#1#2#3#4{\line{\consetup{#1}%
  #3\ %
  \rm\leaders\hbox to .5em{\hss.\hss}\hfil\ \s{#2}\hbox to 3em{\hss#4}}}
% A contents entry's leading material by group depth: a top-level @** section
% (depth -1) sets its title in bold, as cweave does; @* (0) is flush, and deeper
% groups (@*1, @*2, ...) are indented. The \bf is reset to \rm before the page
% number back in \Tline.
\def\consetup#1{\ifnum#1<0 \bf \else
  \ifcase#1 \or\hskip2em\or\hskip4em\or\hskip6em\else\hskip8em\fi \fi}
% \topofcontents and \botofcontents open and close the contents page; redefine
% either in limbo, as in cwebmac, to dress it up (a subtitle, a version line, a
% cover note). The defaults reproduce \con's own former top and bottom.
\def\topofcontents{\centerline{\titlefont\Gtitle}\vskip .7in \vfill}
\def\botofcontents{\vfill}
\def\con{\ifsecs\par\endgroup\vfil\eject\fi % finish the section names
  \iftocopen\immediate\closeout\toc\fi
  \ifcon\iftocopen
    \titletrue \pageno=\contentspagenumber
    \parindent=0pt
    \topofcontents
    \line{\hfil{\smallfont\sectionword}\hbox to 3em{\hss\smallfont\pageword}}%
    \begingroup \rm
      \readcontents\relax
    \endgroup
    \botofcontents
  \fi\fi
  \end}
